% Source: https://artofproblemsolving.com/community/c3h2579477_a_quicker_solution
% Stable identity: aops-topic-2579477-q01
Let $a>b\ge1$ with $a^2-b^2=(a-b)(a+b)=N$, $N$ even.

\emph{The difference cannot be $1$:} if $a-b=1$, then $N=a+b=2b+1$ is odd, a contradiction.

\emph{Difference $2$:} write $N=4m$. Then $a=m+1$, $b=m-1$ give $a^2-b^2=4m=N$, and $b\ge1$ exactly when $m\ge2$, i.e. $N>4$.

\textbf{(a)} $N=400$, $m=100$: $101^2-99^2=400$, so the least difference is $\boxed{2}$.

\textbf{(b)} By the two observations, the least difference is $2$ for every $N=4m$ with $m\ge2$; this is the student's argument in general form, with $x=m-1$, $y=m$, $z=m+1$. For $N=4$ there are no two positive integers at all: $(a-b)(a+b)=4$ with factors of the same parity forces $a-b=a+b=2$, so $b=0$.
